% Part: normal-modal-logic
% Chapter: sequent-calculus
% Section: introduction

\documentclass[../../../include/open-logic-section]{subfiles}

\begin{document}

\olfileid{nml}{seq}{int}

\olsection{引言}

通过添加处理~\iftag{prvBox}{$\Box$\iftag{prvDiamond}{ 和~}{}}{}%\iftag{prvDiamond}{$\Diamond$}{} 的规则，可以扩展命题逻辑的相继式演算。例如，对于~$\Log{K}$，我们在~$\Log{LK}$ 之外加上如下规则：
\[
  \iftag{prvBox}{\iftag{prvDiamond}{% <> and [] primitive
    \Axiom$\Gamma \fCenter \Delta, !A$
    \RightLabel{$\Box$}
    \UnaryInf$\Box\Gamma \fCenter \Diamond\Delta, \Box!A$
    \DisplayProof
    \qquad
    \Axiom$!A, \Gamma \fCenter \Delta$
    \RightLabel{$\Diamond$}
    \UnaryInf$\Diamond!A, \Box\Gamma \fCenter \Diamond\Delta$
    \DisplayProof
}{% only Box primitive
\Axiom$\Gamma \fCenter !A$
\RightLabel{$\Box$}
\UnaryInf$\Box\Gamma \fCenter \Box!A$
\DisplayProof
}}{% only <> primitive
  \Axiom$!A \fCenter \Delta$
  \RightLabel{$\Diamond$}
  \UnaryInf$\Diamond!A \fCenter \Diamond\Delta$
  \DisplayProof
}
\]
对于~$\Log{K}$ 的扩展逻辑，也必须添加额外的规则。

然而，并非每个模态逻辑都有这样的相继式演算。即使是~$\Log{S5}$——它在语义上非常简单（无需使用可达关系即可定义）——目前也尚不知是否存在一种由~$\Log{LK}$ 得到且在无~\Cut 规则时仍然完备的相继式演算。不过，它确实具有一种无切割且完备的\emph{超相继式}演算。

\end{document}
